\chapter{Các công trình liên quan}
\label{ch:related}

Chương này khảo sát các hướng nghiên cứu liên quan tới khóa luận. Trước hết, chúng tôi điểm lại
các mô hình hỏi đáp trên đồ thị tri thức thời gian theo lối truyền thống. Tiếp theo, chúng tôi xem
xét hai hướng khai thác mô hình ngôn ngữ lớn cho suy luận: kết hợp tri thức ngoài và kết hợp bộ
nhớ. Sau đó, chúng tôi giới thiệu các bộ dữ liệu TKGQA tiêu biểu và chỉ ra khoảng trống về dữ liệu
tiếng Việt. Cuối cùng, Mục~\ref{sec:positioning} định vị đóng góp của khóa luận trong bức tranh
nghiên cứu chung.

\section{Các mô hình TKGQA truyền thống}
\label{sec:trad-tkgqa}

Các phương pháp TKGQA truyền thống có thể chia thành hai nhóm chính~\cite{chen2024ari}: nhóm dựa
trên \textit{phân rã câu hỏi} và nhóm dựa trên \textit{biểu diễn nhúng}.

\paragraph{Nhóm phân rã câu hỏi.} Tiêu biểu là TEQUILA~\cite{jia2018tequila}. Phương pháp này phân
rã câu hỏi ban đầu thành các câu hỏi con cùng các ràng buộc thời gian, sử dụng các mô hình KGQA
chuẩn để giải từng câu hỏi con, rồi thực hiện phân tích so sánh để chọn ra đáp án phù hợp nhất.
Cách tiếp cận này phụ thuộc nhiều vào các luật được thiết kế thủ công để phân rã câu hỏi và xử lý
ràng buộc thời gian, nên khó mở rộng sang các dạng câu hỏi đa dạng trong thực tế.

\paragraph{Nhóm biểu diễn nhúng.} Hướng này biểu diễn thực thể, quan hệ và thời gian thành các
vector nhúng, rồi cho điểm các đáp án ứng viên dựa trên độ tương đồng ngữ nghĩa. EmbedKGQA~
\cite{saxena2020embedkgqa} ban đầu được thiết kế cho đồ thị tri thức \emph{tĩnh}: nó tận dụng biểu
diễn nhúng của đồ thị để cải thiện hỏi đáp đa bước, nhưng bỏ qua thông tin thời gian.
CronKGQA~\cite{saxena2021cronkgqa} mở rộng ý tưởng này cho đồ thị tri thức thời gian bằng cách dùng
biểu diễn nhúng thời gian để đánh giá độ tương đồng khi xác định đáp án; phương pháp này hoạt động
tốt với các câu hỏi đơn giản nhưng suy giảm rõ rệt với câu hỏi phức tạp đòi hỏi suy luận thời gian
chuyên biệt. TempoQR~\cite{mavromatis2021tempoqr} khắc phục một phần bằng cách bổ sung thông tin
phạm vi thời gian và làm giàu biểu diễn câu hỏi về mặt ngữ nghĩa.

\paragraph{Hỏi đáp đa hạt thời gian.} MultiQA~\cite{chen2023multitq} hướng tới các câu hỏi có nhiều
mức hạt thời gian (năm, tháng, ngày) thông qua một mô-đun tổng hợp thời gian dựa trên Transformer.
Bên cạnh đó, hướng suy luận trên đồ thị con~\cite{chen2022subgraph} khai thác cấu trúc đồ thị con
quanh thực thể trong câu hỏi để định hướng suy luận. Đáng chú ý, ý tưởng trích xuất và thao tác
trên đồ thị con này cũng là nền tảng cho pha tương tác dựa trên tri thức trong khung ARI
(Mục~\ref{sec:knowledge-based}).

\paragraph{Mô hình ngôn ngữ tiền huấn luyện làm chuẩn so sánh.} Ngoài các mô hình chuyên biệt cho
TKGQA, các mô hình ngôn ngữ tiền huấn luyện như BERT~\cite{devlin2019bert} cũng thường được dùng
làm chuẩn so sánh: biểu diễn câu hỏi sinh từ mô hình được ghép với biểu diễn nhúng của thực thể và
thời gian, rồi chấm điểm với toàn bộ thực thể và mốc thời gian thông qua tích vô hướng. Tuy nhiên,
do không được thiết kế riêng cho suy luận thời gian, các mô hình này thường cho kết quả thấp, đặc
biệt với câu hỏi có đáp án là mốc thời gian.

Nhìn chung, các phương pháp truyền thống đều phụ thuộc vào luật thủ công hoặc biểu diễn nhúng được
huấn luyện trên dữ liệu chuyên biệt. Chúng yêu cầu giai đoạn huấn luyện tốn kém, khó thích nghi với
miền tri thức mới, và đặc biệt còn gặp khó khăn với các câu hỏi đòi hỏi suy luận thời gian tinh vi
nhiều bước~\cite{chen2022subgraph}. Đây là động lực để chuyển sang khai thác năng lực suy luận linh
hoạt của mô hình ngôn ngữ lớn.

\section{Suy luận của mô hình ngôn ngữ lớn với tri thức ngoài}
\label{sec:llm-external}

Để giảm ảo giác và bổ sung tri thức, nhiều nghiên cứu tích hợp tri thức ngoài vào mô hình ngôn ngữ
lớn. Theo cách phân loại của~\cite{baek2023kaping}, có hai hình thức: tiêm tri thức \textit{tường
minh} và tiêm tri thức \textit{ngầm}.

\paragraph{Tiêm tri thức tường minh.} Tri thức liên quan được cung cấp trực tiếp cho mô hình thông
qua câu lệnh. KAPING~\cite{baek2023kaping} truy hồi các bộ ba liên quan tới câu hỏi từ đồ thị tri
thức rồi nối vào câu lệnh, cho phép hỏi đáp ở chế độ \textit{zero-shot} mà không cần huấn luyện.
CoK (Chain-of-Knowledge)~\cite{li2023cok} trước hết đánh giá độ tin cậy của câu trả lời, và khi cần
thì dùng mô hình ngôn ngữ lớn phân rã câu hỏi rồi sinh các truy vấn để trích xuất thông tin từ
nguồn tri thức ngoài. ChatKBQA~\cite{luo2023chatkbqa} tinh chỉnh mô hình ngôn ngữ lớn để sinh ra
các truy vấn logic có thể thực thi trực tiếp trên đồ thị tri thức nhằm thu được đáp án.

\paragraph{Khai thác tương tác và tri thức biểu trưng.} ToG (Think-on-Graph)~\cite{sun2023tog} coi
mô hình ngôn ngữ lớn như một tác tử, cho phép nó khám phá tương tác các thực thể và quan hệ liên
quan trên đồ thị tri thức rồi suy luận dựa trên tri thức truy xuất được. Symbol-LLM~
\cite{xu2023symbolllm} đề xuất khung tinh chỉnh hai giai đoạn nhằm tích hợp tri thức biểu trưng
(\textit{symbolic}) vào mô hình ngôn ngữ lớn, qua đó tăng cường năng lực suy luận.

\paragraph{Hạn chế.} Việc tích hợp tri thức tuy giảm được ảo giác nhưng không phải không có thách
thức. Tiêm tri thức tường minh thường khó thu được thông tin liên quan, chất lượng cao và bị giới
hạn bởi độ dài ngữ cảnh hữu hạn; trong khi tiêm tri thức ngầm thường đòi hỏi tinh chỉnh tham số,
vốn tốn kém. Với đồ thị tri thức thời gian, một câu hỏi có thể liên quan tới hàng nghìn sự kiện,
khiến việc nhồi tri thức vào câu lệnh trở nên bất khả thi và dễ gây nhiễu. Khung ARI khắc phục bằng
cách \emph{tách bạch} suy luận thành hai phần -- phần liên quan tri thức và phần không phụ thuộc tri
thức -- nhờ đó tránh được đồng thời cả hai hạn chế trên.

\section{Suy luận của mô hình ngôn ngữ lớn với bộ nhớ}
\label{sec:llm-memory}

Vì mô hình ngôn ngữ lớn thiếu bộ nhớ dài hạn, nhiều nghiên cứu trang bị cho nó cơ chế bộ nhớ để
lưu và tái sử dụng kinh nghiệm. MemoChat~\cite{lu2023memochat} xây dựng và cập nhật một bản ghi nhớ
(\textit{memo}) có cấu trúc, phân loại các đoạn hội thoại trong quá khứ rồi truy xuất theo chủ đề
và tóm tắt. Reflexion~\cite{shinn2023reflexion} khai thác một bộ nhớ làm việc để lưu kinh nghiệm
cho một tác vụ cụ thể, cải thiện hiệu năng qua nhiều lần thử; tuy nhiên lịch sử lưu trong bộ nhớ
làm việc khó mang lại lợi ích cho các tác vụ có mục tiêu khác nhau. MemPrompt~\cite{madaan2022memprompt}
thiết kế một bộ nhớ bền vững lưu phản hồi của người dùng để nhắc lại cho mô hình và cải thiện liên
tục. RLEM~\cite{zhang2023rlem} sử dụng một bộ nhớ kinh nghiệm gắn với môi trường, hỗ trợ ra quyết
định trong tương lai kể cả với mục tiêu tác vụ khác nhau. Thought Propagation~\cite{yu2023tp} nhấn
mạnh khả năng khám phá và áp dụng tri thức từ các lời giải tương tự.

\paragraph{Hạn chế.} Các phương pháp tăng cường bộ nhớ hiện nay phần lớn giới hạn ở việc
\emph{tiếp nhận thụ động} thông tin lịch sử, bỏ qua việc \emph{chủ động kiến tạo} tri thức trừu
tượng từ kinh nghiệm trước đó. Xuất phát từ thuyết kiến tạo (Mục~\ref{sec:memory-bg}), khung ARI
cung cấp cho mô hình một quá trình học chủ động và liên tục, sản sinh tri thức ở dạng trừu tượng và
tổng quát hóa được -- đây chính là điểm khác biệt cốt lõi mà khóa luận kế thừa và phát triển cho
tiếng Việt.

\section{Các bộ dữ liệu TKGQA}
\label{sec:datasets}

Hai bộ dữ liệu TKGQA được sử dụng rộng rãi nhất là CRONQUESTIONS và MULTITQ.

\textbf{CRONQUESTIONS}~\cite{saxena2021cronkgqa} là bộ dữ liệu hỏi đáp trên đồ thị tri thức thời
gian xây dựng từ Wikidata~\cite{vrandecic2014wikidata}, trong đó các thực thể và mốc thời gian xuất
hiện trong câu hỏi đều được chú thích. Bộ dữ liệu gồm cả câu hỏi đơn giản lẫn phức tạp, được tổ
chức thành năm loại: thực thể đơn, thời gian đơn, đồng thời (\textit{time join}), đầu/cuối
(\textit{first/last}) và trước/sau (\textit{before/after}).

\textbf{MULTITQ}~\cite{chen2023multitq} là bộ dữ liệu câu hỏi thời gian phức tạp với thông tin thời
gian \emph{đa hạt} (năm, tháng, ngày), xây dựng từ kho sự kiện ICEWS~\cite{boschee2015icews}. So
với các bộ dữ liệu trước, MULTITQ nổi bật ở quy mô lớn, hệ quan hệ phong phú và nhiều mức hạt thời
gian, nhờ đó phản ánh tốt hơn các tình huống thực tế.

\paragraph{Xây dựng và lọc đồ thị tri thức thời gian.} Chất lượng của một bộ dữ liệu TKGQA phụ
thuộc trực tiếp vào chất lượng đồ thị tri thức nền. Trong lĩnh vực \textit{hoàn thiện đồ thị tri
thức thời gian} (\textit{temporal knowledge base completion}), các nghiên cứu như~\cite{lacroix2020tkbc}
thường áp dụng bước tiền xử lý lọc bỏ các quan hệ và thực thể quá hiếm để giữ lại phần đồ thị dày
đặc, giàu thông tin và ổn định về mặt thống kê. Khóa luận kế thừa tinh thần lọc này khi trích xuất
đồ thị tri thức thời gian tiếng Việt từ Wikidata: chỉ giữ các sự kiện có định tố thời gian, đồng
thời loại bỏ quan hệ và thực thể có tần suất thấp (chi tiết ở Mục~\ref{sec:buildkg}).

Điểm chung của hai bộ dữ liệu trên là đều dành cho \textbf{tiếng Anh}, với nhãn thực thể và quan hệ
bằng tiếng Anh, xây dựng từ Wikidata hoặc ICEWS. Cho tới nay \emph{chưa có} bộ dữ liệu TKGQA nào
dành cho tiếng Việt. Đây chính là khoảng trống mà khóa luận lấp đầy bằng bộ dữ liệu CronQ-VN
(Mục~\ref{sec:cronqvn}), được trích xuất từ Wikidata với quy trình lọc lấy cảm hứng từ kỹ thuật
hoàn thiện đồ thị tri thức thời gian~\cite{lacroix2020tkbc}. Bảng~\ref{tab:datasets} so sánh các
bộ dữ liệu.

\begin{table}[ht]
\centering
\caption{So sánh các bộ dữ liệu TKGQA.}
\label{tab:datasets}
\begin{tabular}{L{3.2cm} C{2.6cm} C{2.8cm} C{2cm} C{2.4cm}}
\toprule
\textbf{Bộ dữ liệu} & \textbf{Nguồn KG} & \textbf{Hạt thời gian} & \textbf{Ngôn ngữ} & \textbf{\# câu hỏi} \\
\midrule
CRONQUESTIONS & Wikidata & Năm & Anh & $\sim$410K \\
MULTITQ & ICEWS & Năm/Tháng/Ngày & Anh & $\sim$500K \\
\textbf{CronQ-VN} & Wikidata & Năm & \textbf{Việt} & \todo{số liệu} \\
\bottomrule
\end{tabular}
\end{table}

\section{Định vị đề tài}
\label{sec:positioning}

Bảng~\ref{tab:positioning} so sánh khóa luận với một số phương pháp đại diện trên các tiêu chí then
chốt. So với các mô hình truyền thống, khóa luận khai thác mô hình ngôn ngữ lớn nên không cần huấn
luyện lại biểu diễn nhúng cho từng miền. So với các phương pháp dùng mô hình ngôn ngữ lớn khác,
khóa luận kế thừa khung ARI -- vốn kết hợp cả tương tác tri thức lẫn bộ nhớ kinh nghiệm trừu tượng.

\begin{table}[ht]
\centering
\caption{Định vị khóa luận so với các phương pháp đại diện.}
\label{tab:positioning}
\begin{tabular}{L{3.4cm} C{2.6cm} C{1.8cm} C{1.8cm} C{2.2cm}}
\toprule
\textbf{Phương pháp} & \textbf{Nhóm} & \textbf{Dùng LLM} & \textbf{Bộ nhớ} & \textbf{Cần huấn luyện} \\
\midrule
CronKGQA & Nhúng & Không & Không & Có \\
TempoQR & Nhúng & Không & Không & Có \\
KG-RAG & LLM + tri thức & Có & Không & Không \\
ReAct-KB & LLM tác tử & Có & Không & Không \\
ARI (gốc) & LLM + bộ nhớ & Có & Có & Không \\
\textbf{Khóa luận} & \textbf{LLM + bộ nhớ} & \textbf{Có (cục bộ)} & \textbf{Có} & \textbf{Không$^{*}$} \\
\bottomrule
\end{tabular}
\end{table}

\noindent{\footnotesize $^{*}$ Khung suy luận không cần huấn luyện; riêng bộ lọc cross-encoder đề
xuất (đóng góp thứ ba) có một bước huấn luyện nhẹ cho bộ chấm điểm.}

\medskip

Trên cơ sở khảo sát trên, khóa luận xác định ba khoảng trống và đề ra hướng giải quyết tương ứng:

\begin{enumerate}
  \item \textbf{Khoảng trống về ngôn ngữ.} Toàn bộ các bộ dữ liệu và hệ thống TKGQA hiện có đều cho
  tiếng Anh. Khóa luận xây dựng bộ dữ liệu TKGQA tiếng Việt CronQ-VN và một hệ thống hỏi đáp tương
  ứng.
  \item \textbf{Khoảng trống về triển khai cục bộ.} Các hệ thống dựa trên mô hình ngôn ngữ lớn phổ
  biến phụ thuộc vào API thương mại. Khóa luận vận hành toàn bộ khung ARI bằng mô hình ngôn ngữ
  chạy cục bộ, kèm cơ chế liên kết thực thể tiếng Việt cho câu hỏi mở.
  \item \textbf{Khoảng trống về lọc hành động nhận biết ngữ cảnh.} Bộ lọc dựa trên độ tương đồng
  cosine trong ARI gốc không nhận biết bước suy luận hiện tại. Khóa luận đề xuất thay thế bằng bộ
  lọc dựa trên cross-encoder, mã hóa đồng thời hành động và lịch sử suy luận.
\end{enumerate}

Ba khoảng trống này định hình ba đóng góp của khóa luận, được trình bày chi tiết trong
Chương~\ref{ch:method}.
